% app_b_t1.tex -- Appendix B: T1 audit details (tables, ceiling construction, proofs).
% Content moved verbatim from 03_t1_proxy.tex.

\section{The Checkpoint-Noise Proxy: Full Tables and Proofs}
\label{app:t1}

\subsection{Home replication table}
\label{app:t1-home}

% TABLE-t1a: home replication. Source: tables/t1a_sn_replication_summary.csv
\begin{table}[!htbp]
\centering
% TABLE-t1a: S&N home replication (source tables/t1a_sn_replication_summary.csv; per-task
% breakdown in appendix from t1a_sn_replication_per_task.csv)
\caption{\textbf{Home-setting replication of the checkpoint-noise proxy.} Cross-task correlation
between the final-checkpoint proxy $x$ and true seed noise $y$ on Signal \& Noise's own 1B runs,
by arm and metric space, plus the calibration ratio $y/x$ (median $0.98$). The original claim
reproduces.}
\label{tab:t1a}
\footnotesize
\begin{tabular}{@{}llccccc@{}}
\toprule
arm & metric set & $n$ tasks & $R$ (raw) & $R^2$ (raw) & $R$ (log) & $R^2$ (log) \\
\midrule
init-seed (10 runs)   & primary-like  & 9  & 0.996 & 0.992 & 0.902 & 0.814 \\
init-seed (10 runs)   & bits-per-byte & 18 & 0.953 & 0.908 & 0.762 & 0.581 \\
data-order (9 runs)   & primary-like  & 9  & 0.991 & 0.981 & 0.939 & 0.882 \\
data-order (9 runs)   & bits-per-byte & 18 & 0.951 & 0.904 & 0.779 & 0.607 \\
\bottomrule
\end{tabular}
\end{table}

For each of the 10 init-seed runs, $x$ is the relative SD of the final 30 checkpoints (500-step
grid); $y$ is the across-run relative SD of final scores within the arm; the statistic is the
cross-task Pearson correlation, in raw and log space.

\subsection{Porting to DataDecide: per-size table}
\label{app:t1-port}

Porting protocol: for each of 14 sizes $\times$ 25 recipes, $x$ is the relative SD
of the default seed's last $n$ grid checkpoints ($n = \min(5, \lfloor \mathrm{grid}/2 \rfloor)
\ge 3$), $y$ is the 3-seed relative SD of final scores, and $R$ is computed across the 10 OLMES
tasks per cell (350 cells).

% TABLE-t1b: ported correlation + MC ceiling. Source: tables/t1b_summary_by_size.csv
\begin{table}[!htbp]
\centering
% TABLE-t1b: per-size R_log median, MC ceiling, de-attenuated lambda, calibration ratio
% (sources: tables/t1b_summary_by_size.csv; lambda from review-verify/w5_lambda_v2.py;
% bias-corrected y/x from review-verify/batch2_remaining.py, median-basis)
\caption{\textbf{Porting the proxy to DataDecide.} Per size: median cross-task
$R_{\log}$ over 25 recipes; the Monte-Carlo ceiling on $R$ for a \emph{perfect} proxy at 3
seeds; the de-attenuated $\hat\lambda = \hat R/(\text{ceiling median})$ with recipe-cluster
bootstrap 95\% interval ($^*$capped at $1.00$: an uncapped value above $1$ signals that the
ceiling understates the attenuation, consistent with the porting protocol's shared-seed
dependence, Appendix~\ref{app:t1-lambda}); and the calibration ratio $y/x$, raw and with the
median-basis log-space correction. Across sizes, $60$--$96\%$ of cells are compatible with a
perfect proxy and $16$--$48\%$ indistinguishable from null. No size's interval excludes or
certifies the advertised $0.9$.}
\label{tab:t1b}
\footnotesize
\setlength{\tabcolsep}{3.5pt}
\begin{tabular}{@{}lccccc@{}}
\toprule
size & med.\ $R_{\log}$ & ceiling med.\ & $\hat\lambda$ [95\% CI] & med.\ $y/x$ & $y/x$ (corr.) \\
\midrule
4M & 0.74 & 0.80 & 0.92 [0.79, 1.04] & 1.60 & 1.60 \\
6M & 0.54 & 0.79 & 0.69 [0.57, 0.90] & 2.21 & 2.21 \\
8M & 0.62 & 0.79 & 0.79 [0.54, 0.94] & 1.20 & 1.32 \\
10M & 0.76 & 0.82 & 0.92 [0.69, 1.03] & 1.45 & 1.59 \\
14M & 0.70 & 0.82 & 0.86 [0.65, 0.97] & 1.08 & 1.18 \\
16M & 0.75 & 0.79 & 0.95 [0.82, 1.00] & 1.53 & 1.69 \\
20M & 0.74 & 0.78 & 0.94 [0.80, 1.06] & 1.42 & 1.56 \\
60M & 0.77 & 0.75 & 1.00$^*$ [0.89, 1.11] & 2.05 & 2.26 \\
90M & 0.66 & 0.71 & 0.93 [0.75, 1.08] & 1.70 & 1.87 \\
150M & 0.68 & 0.73 & 0.93 [0.80, 1.10] & 1.99 & 2.19 \\
300M & 0.68 & 0.72 & 0.95 [0.76, 1.09] & 2.17 & 2.38 \\
530M & 0.72 & 0.72 & 1.00 [0.86, 1.09] & 1.36 & 1.49 \\
750M & 0.74 & 0.75 & 0.98 [0.83, 1.07] & 0.95 & 1.04 \\
1B & 0.53 & 0.66 & 0.80 [0.42, 0.96] & 1.54 & 1.69 \\
\bottomrule
\end{tabular}
\end{table}

The Monte-Carlo ceiling is the distribution of $R$ one would observe if the proxy were
\emph{perfect}, given only the sampling error of a 3-seed SD ($\chi^2$ with 2 degrees of
freedom) and an $n$-checkpoint SD.

\subsection{The attenuation model and proofs}
\label{app:t1-proof}

The ceiling is a theorem, and its statement is short; Proposition~\ref{prop:t1-bound} and
Corollary~\ref{cor:t1-n3} appear in the main text (Section~\ref{sec:t1-port}). The model behind
them: write the log proxy SD and log seed SD per task as $U_i = a_i + \varepsilon_i$ and
$Z_i = b_i + \eta_i$, where $a_i = \log \tau_i$ is the
task's latent noise level with cross-task variance $\sigma_a^2$, and the estimation errors are
$\varepsilon_i = \tfrac12\log(\chi^2_{\nu_x}/\nu_x)$, $\eta_i = \tfrac12\log(\chi^2_{\nu_y}/\nu_y)$
with $\nu_x = m-1$ ($m$ checkpoints) and $\nu_y = n-1$ ($n$ seeds). Then
$\mathrm{Var}(\varepsilon) = \psi_1(\nu_x/2)/4$ exactly ($\psi_1$ the trigamma function; this is
the log-space analogue of the $c_4$ small-sample SD bias, and it is larger:
$b(2) = -0.289$ nat vs.\ $\log c_4(3) = -0.121$). This is an errors-in-variables model with
\emph{both} arms noisy, and the population correlation is exactly attenuated. The Fisher-$z$
sample-complexity statement is an asymptotic, model-based approximation at finite $k$, since the
observed log-SD pairs are only approximately bivariate normal; the bound is stated against the
most favorable alternative ($\lambda = 1$, $A$ known), an oracle that only helps the verifier.
In that limit experiment the power of any level-$\alpha$ certificate of $\lambda \ge R_0$ from
$k$ tasks is at most
$\Phi\!\bigl(\sqrt{k-3}\,[\operatorname{atanh} A - \operatorname{atanh}(A R_0)] -
z_{1-\alpha}\bigr)$; inversion gives the sample-complexity bound of
Proposition~\ref{prop:t1-bound}. (With unequal latent dispersions, replace $\sigma_a^2$ by the
per-arm $\sigma_a^2, \sigma_b^2$ in each factor of $A$.)

\begin{proof}[Proof of Proposition~\ref{prop:t1-bound} (sketch)]
$\log S = \log\sigma + \tfrac12\log(\chi^2_\nu/\nu)$ gives $\mathrm{Var}(\log S) = \psi_1(\nu/2)/4$.
The pairs $(U_i, Z_i)$ are two noisy views of latents correlated at $\lambda$, hence
$\rho = \lambda\sigma_a^2/\sqrt{(\sigma_a^2 + \mathrm{Var}\,\varepsilon)(\sigma_a^2 +
\mathrm{Var}\,\eta)} = \lambda A$: errors-in-variables attenuation with both arms noisy. Under Fisher's transform
$\operatorname{atanh}\hat r \to_d \mathcal N(\operatorname{atanh}(\lambda A), (k-3)^{-1})$; the
composite null $\lambda \le R_0$ is least favorable at $\lambda = R_0$, and the $z$-test is UMP in
the limit experiment, giving the power bound; inversion gives the bound on $k$. Treating $A$ as
known is favorable to the verifier (estimating $\sigma_a$ from the same data can only degrade
power), so the bound is a valid lower bound on the true requirement.
\end{proof}

The full arithmetic behind Corollary~\ref{cor:t1-n3}: on the DataDecide panel ($k = 10$ OLMES
tasks, $m \le 5$ checkpoints, bias-corrected median
cross-task dispersion $\hat\sigma_a = 0.74$ in log units), a 3-seed arm certifying $\lambda \ge 0.9$
at $\alpha = 0.05$ has power $0.089$ under this asymptotic model even if the proxy is perfect;
certification needs
$k = 266$--$1{,}516$ tasks by size (median $503$; most favorable decile $173$). Fixing $k = 10$:
with $m = 5$ no seed budget suffices for any panel with $\sigma_a \le 2.04$; with $m = 30$ one needs
$n \ge 135$ seeds per arm at $\sigma_a = 0.75$ ($n \ge 27$ at $\sigma_a = 1.0$). The home budget of
Signal \& Noise itself ($n = 10$, $m = 30$, $k = 9$--$10$) attains power $0.31$--$0.70$ for
$\sigma_a \in [0.75, 1.25]$: the advertised $R \ge 0.9$ was a point estimate, not a certified bound,
even in the original setting. For $99.7\%$ of our 350 cells the perfect proxy's \emph{population}
correlation is already below $0.9$. The closed form reproduces the Monte-Carlo ceiling
($0.66$--$0.82$) to within $0.03$.

\paragraph{Finite-sample calibration of the certificate.}
\label{app:t1-mc}
The limit experiment is an asymptotic idealization; we calibrate the certificate's actual
operating characteristics at the panel's own operating point by Monte-Carlo ($400{,}000$
replicates per cell, simulating the exact model of Proposition~\ref{prop:t1-bound} at $k=10$
tasks, $m=5$ checkpoints, $n=3$ seeds, $\sigma_a = 0.74$, $R_0 = 0.9$, $\alpha = 0.05$). At the
boundary $\lambda = R_0$ the one-sided Fisher-$z$ certificate's actual Type-I error is $0.069$
against the nominal $0.05$ (the small-$k$ approximation is mildly liberal); its power against a
perfect proxy ($\lambda = 1$) is $0.124$, close to the nominal $0.089$. Under equicorrelated
task noise ($\rho_{\mathrm{task}} = 0.2$, $0.5$) the test turns conservative (Type-I
$0.046$, $0.019$) and power degrades further ($0.083$, $0.036$); at $k = 50$ the nominal power
is still only $0.188$. The certificate is thus uninformative in every regime we can simulate,
and the asymptotic bound of Proposition~\ref{prop:t1-bound} is, if anything, optimistic about
the verifier. (Script: \texttt{verify-scripts/r4\_prop1\_mc\_calibration.py} in the released
reproduction package.)

\subsection{De-attenuation, the shared-seed caveat, and the leave-one-seed-out check}
\label{app:t1-lambda}

The $\hat\lambda$ of Table~\ref{tab:t1b} divides the observed median $R_{\log}$ by the
Monte-Carlo ceiling median, i.e.\ by what a perfect proxy would typically read at this seed
budget. Two estimators bracket the same quantity: a model-based version, evaluating the $A$ of
Proposition~\ref{prop:t1-bound} per size with per-arm bias-corrected dispersions, gives a
median of $1.03$ across sizes; the ceiling ratio gives $0.93$. The gap is the visible symptom
of a dependence in the porting protocol itself: the proxy $x$ is computed on the
\emph{default} seed's run, whose final checkpoint also enters the seed SD $y$, so the two
arms' estimation errors are positively dependent, the observed $R$ is inflated, and both
de-attenuations are optimistic (Proposition~\ref{prop:t1-bound} assumes independent arms, so
its ceiling is an upper bound here). Under a leave-one-seed-out variant (the proxy computed on
seed $s$, the seed SD on the other two seeds, averaged over the three choices; the ceiling
drops with $\nu_y = 1$), $\hat\lambda$ has median $0.89$ across sizes, and at the smallest
scales the intervals fall short of $0.9$ (4M [$0.70$, $0.86$]; 6M [$0.45$, $0.70$]; 10M
[$0.59$, $0.88$]; Table~\ref{tab:lambda-loo}). We therefore read Table~\ref{tab:t1b}'s
$\hat\lambda$ as consistent with the advertised level, not as evidence for it; certification
remains impossible either way (Corollary~\ref{cor:t1-n3}).

% TABLE-lambda-loo: source review-verify/w5_loo_lambda.py (B=2000 recipe bootstrap; MC ceiling
% nsim=200 per cell under the LOO construction).
\begin{table}[!htbp]
\centering
\caption{\textbf{The leave-one-seed-out porting.} Per size: median cross-task $R_{\log}$ with
proxy and seed arms sharing no run; the corresponding perfect-proxy ceiling ($\nu_y = 1$); and
the de-attenuated $\hat\lambda$ with recipe-cluster bootstrap interval. Point estimates above
$1$ at 90M/150M reflect the noisier 2-seed ceiling, not super-perfect proxies.}
\label{tab:lambda-loo}
\footnotesize
\setlength{\tabcolsep}{4pt}
\begin{tabular}{@{}lccc@{}}
\toprule
size & med.\ $R_{\log}$ (LOO) & ceiling med.\ (LOO) & $\hat\lambda$ [95\% CI] \\
\midrule
4M & 0.54 & 0.67 & 0.79 [0.70, 0.86] \\
6M & 0.32 & 0.65 & 0.49 [0.45, 0.70] \\
8M & 0.51 & 0.61 & 0.83 [0.60, 0.92] \\
10M & 0.50 & 0.66 & 0.76 [0.59, 0.88] \\
14M & 0.58 & 0.65 & 0.89 [0.70, 1.00] \\
16M & 0.52 & 0.61 & 0.86 [0.61, 0.95] \\
20M & 0.58 & 0.61 & 0.95 [0.86, 1.06] \\
60M & 0.55 & 0.57 & 0.96 [0.88, 1.17] \\
90M & 0.55 & 0.53 & 1.03 [0.98, 1.16] \\
150M & 0.59 & 0.55 & 1.07 [0.83, 1.16] \\
300M & 0.47 & 0.54 & 0.88 [0.66, 1.08] \\
530M & 0.52 & 0.53 & 0.97 [0.72, 1.14] \\
750M & 0.55 & 0.58 & 0.96 [0.69, 1.10] \\
1B & 0.31 & 0.47 & 0.66 [0.51, 0.94] \\
\bottomrule
\end{tabular}
\end{table}

\paragraph{Bias-correction convention.} For median quantities we apply the median of the
log-$\chi^2$ law (the median of $\tfrac12\log(\chi^2_\nu/\nu)$ is $-0.183$ at $\nu = 2$ and
$-0.088$ at $\nu = 4$; the $y/x$ correction factor is $1.10$ at $m = 5$, $1.07$ at $m = 4$,
and $1$ at $m = 3$ where the two arms' biases cancel), not its mean; mean-based corrections
apply only to mean quantities, and the variance-scale estimators of Appendix~\ref{app:t3}
need none.

\subsection{Continuous metrics: trend contamination}
\label{app:t1-ppl}

On DataDecide's 11-domain perplexity tables (log space), the raw proxy fails in the
\emph{opposite} direction: the median $y/x$ is $0.14$--$0.89$ across sizes, i.e., the proxy
\emph{overestimates} seed noise by $2$--$7\times$. The cause is mechanical: on coarse
evaluation grids the final checkpoints are still improving, so the checkpoint spread mixes seed
noise with leftover learning trend. Detrending (residual SD after a linear fit on the final
window) repairs the calibration at large scale (median ratio $1.03$--$1.08$ at 750M/1B) but
not at $\le 530$M, where ratios still scatter $0.54$--$2.8$ and only $26$--$64\%$ of cells fall
within $2\times$. The direction of the proxy's bias is therefore metric- and grid-dependent,
not a stable property one could correct for.

The repair is also non-monotone in scale: detrended ratios at 60M/300M ($2.82$/$2.26$) exceed
those at 90M/150M ($1.22$/$1.52$). This is not a grid-density effect (the window has five
checkpoints at both 60M and 90M), and it persists when linear detrending is replaced by a
trend-invariant first-difference proxy ($2.41$/$2.67$ at 60M/300M), so it is not an artifact
of the linear fit either. We have no mechanism-level account; the license conditions of
Section~\ref{sec:t1-license} are therefore stated at the median level only.

% TABLE-t1c: ppl raw vs detrended. Source: tables/t1c_ppl_summary_by_size.csv
\begin{table}[!htbp]
\centering
% TABLE-t1c: per-size ppl proxy ratios, raw vs detrended (source
% tables/t1c_ppl_summary_by_size.csv)
\caption{\textbf{Trend contamination on perplexity.} Calibration ratio $y/x$ of the checkpoint
proxy on 11-domain log-perplexity, raw and after detrending the final checkpoint window. Raw
ratios of $0.14$--$0.89$ mean the proxy overestimates seed noise $2$--$7\times$; detrending
restores $1.03$--$1.08$ only at ${\ge}750$M. Ratios are uncorrected; the median-basis
correction of the previous subsection shifts them by ${\approx}+10\%$ at $m = 5$ without
changing any direction.}
\label{tab:t1c}
\footnotesize
\begin{tabular}{@{}lcccc@{}}
\toprule
& \multicolumn{2}{c}{median ratio $y/x$} & \multicolumn{2}{c}{cells within a $2\times$ band} \\
\cmidrule(lr){2-3}\cmidrule(lr){4-5}
size & raw & detrended & raw & detrended \\
\midrule
4M & 0.22 & 0.62 & 17\% & 44\% \\
6M & 0.29 & 1.09 & 23\% & 44\% \\
8M & 0.19 & 0.68 & 18\% & 51\% \\
10M & 0.19 & 0.85 & 18\% & 43\% \\
14M & 0.14 & 0.54 & 12\% & 45\% \\
16M & 0.21 & 0.76 & 21\% & 53\% \\
20M & 0.45 & 1.53 & 45\% & 51\% \\
60M & 0.89 & 2.82 & 64\% & 26\% \\
90M & 0.41 & 1.22 & 38\% & 64\% \\
150M & 0.43 & 1.52 & 34\% & 51\% \\
300M & 0.70 & 2.26 & 45\% & 41\% \\
530M & 0.26 & 1.75 & 23\% & 46\% \\
750M & 0.25 & 1.08 & 21\% & 63\% \\
1B & 0.16 & 1.03 & 14\% & 54\% \\
\bottomrule
\end{tabular}
\end{table}

\subsection{Independent cross-check on PolyPythias}
\label{app:t1-pp}

PolyPythias provides 10-seed ground truth. Comparing each run's final-window step noise ($x$,
from checkpoints past $0.72\times$ of training) against the 10-seed final-score SD ($y$) over 9
available size$\times$task cells, the median ratio is $1.80$ (IQR $1.22$--$2.55$): the proxy
underestimates by $1.8\times$ on an entirely independent suite. The worst cells are saturated
tasks (\textsc{blimp} at 70m/160m reaches $4.2$--$5.2\times$) where step noise collapses once
accuracy plateaus but seed noise does not. The failure is therefore not an artifact of
DataDecide's grid or seed semantics.
